% Sovrapposizione di due onde della stessa frequenza e della stessa ampiezza:
% in fase (interferenza costruttiva) e in opposizione di fase (distruttiva).
\begin{tikzpicture}[>=Stealth, font=\small, x=0.9cm, y=0.75cm]
  \foreach \col/\segno/\titolo in {0/1/{(a) in fase: interferenza costruttiva}, 7.2/-1/{(b) in opposizione di fase: distruttiva}} {
    \begin{scope}[xshift=\col cm]
      \node[anchor=west] at (-0.2,1.3) {\titolo};
      % prima onda
      \draw[richiamo] (0,0) -- (6,0);
      \draw[formula, thick, domain=0:6, samples=100] plot (\x, {0.5*sin(\x*120)});
      \node[formula, right] at (6,0) {onda 1};
      % seconda onda
      \draw[richiamo] (0,-1.5) -- (6,-1.5);
      \draw[medica, thick, domain=0:6, samples=100] plot (\x, {-1.5 + \segno*0.5*sin(\x*120)});
      \node[medica, right] at (6,-1.5) {onda 2};
      \node at (3,-2.35) {$+$};
      % somma
      \draw[richiamo] (0,-3.6) -- (6,-3.6);
      \draw[esempio, very thick, domain=0:6, samples=100] plot (\x, {-3.6 + (1+\segno)*0.5*sin(\x*120)});
      \node[esempio, right] at (6,-3.6) {somma};
    \end{scope}
  }
\end{tikzpicture}
